\chapter{Desiderata}
\label{ch:desiderata}

Before building the calculus, we state what it must achieve.  This
chapter lists eight normative requirements---\emph{desiderata}---that
the design is held accountable to.  Each desideratum is stated as a
named, numbered property so that later chapters can cite it explicitly
when discharging it.  The desiderata are not axioms of the calculus
(those come in Chapter~\ref{ch:world-models}).  They are requirements
\emph{on} the calculus: properties that a well-designed framework for
uncertain reasoning ought to satisfy, and that we claim this one does
satisfy, under conditions that will be made precise.

\paragraph{Why desiderata, not axioms?}
An axiom is a starting point for derivation.  A desideratum is a
target for design.  The distinction matters because several of the
properties below are not simultaneously achievable in full generality:
the no-go theorem (Chapter~\ref{ch:no-go}) shows that local chaining
cannot be universally complete, and the five chaining-limitation
theorems (Chapter~\ref{ch:error-transport}) show that scalar
propagation inevitably loses information.  The desiderata say what
we \emph{want}; the theorems say what we \emph{get} and where we
\emph{stop}.  The interplay between the two is what gives the book
its shape.

\paragraph{Origin.}
These desiderata crystallized during the formalization of Probabilistic
Logic Networks~\cite{goertzel2009pln}.  The original PLN design
implicitly assumed most of them---evidence-based reasoning, the
strength-confidence decomposition, revision as evidence combination,
the idea that inference rules are derived rather than primitive.  What
PLN v0.9 did not do was state these requirements explicitly, prove them
under typed conditions, or delineate where they fail.  The present list
makes the implicit contract explicit.

\paragraph{Independence.}
The eight desiderata are logically independent: for each one, there
exists a system satisfying the other seven while violating it.  These
independence witnesses are noted after each statement.  The
independence claim is not a theorem of the formal development---it is a
design observation that justifies each desideratum's inclusion as
non-redundant.

\bigskip

The desiderata are organized in three layers.  The first two
(\textbf{D1}--\textbf{D2}) define the semantic kernel.  The next two
(\textbf{D3}--\textbf{D4}) constrain the evidence architecture.
The final four (\textbf{D5}--\textbf{D8}) constrain inference,
retraction, conservation, and self-knowledge.

\maplecourt{%
Each desideratum will be illustrated with Maple Court's building
model.  The illustrations are not proofs---they are concrete
instances showing that the desiderata arise naturally in a realistic
evidence-management scenario, not as abstract stipulations.}


%% ═══════════════════════════════════════════════════════════════
\section{D1: State Primacy}
\label{sec:d1}
%% ═══════════════════════════════════════════════════════════════

\begin{definition}[Desideratum D1 --- State Primacy]
\label{def:d1}
A reasoning system should carry a \emph{revisable state} that can be
updated, queried, partially retracted, and explained.  Truth-value
summaries are derived from states, never stored as the primary
representation.
\end{definition}

\noindent
The philosophical motivation is elementary: a bare number has no parts.
You cannot ask \emph{where} a strength of $0.83$ came from, which
observations contributed to it, or what would happen if one source were
retracted.  A state has structure.  It records what was observed, by
whom, and when.  The four verbs of the calculus---revise, query,
forget, explain---all require state; the first two require it
minimally, the last two require it essentially.

\paragraph{What D1 rules out.}
Any system that stores only a scalar summary (a probability, a
strength-confidence pair, a log-odds value) as its primary
representation violates D1.  Classical probability calculus violates D1
when probabilities are treated as fundamental rather than derived from
evidence.  PLN v0.9 satisfies D1 in spirit (it carries evidence
counts) but does not formalize the state operations that D1 enables
(forgetting, explanation).

\paragraph{Discharge.}
D1 is discharged by the \code{WorldModel} typeclass
(Chapter~\ref{ch:world-models}), which takes $\mathrm{State}$ as a
first-class typed parameter with \code{revise}, \code{empty}, and
\code{extract} as operations.

\paragraph{Independence witness.}
A lookup-table system that maps queries directly to real-valued
answers (no state, no revision) satisfies D2--D8 trivially but
violates D1: there is nothing to revise, retract, or explain.

\maplecourt{%
Each Maple Court apartment carries a private state over questions such
as ``Is there a pipe leak?'' and ``Is the shower running?''  The state
is not the strength $0.83$---it is the evidence pair $(5, 1)$: five
leak alarms and one no-alarm reading.  The state can be revised
(a new alarm arrives), queried (what is the current leak evidence?),
forgotten (retract the reading from the miscalibrated sensor), and
explained (the evidence came from sensors 3B-north and 3B-south,
installed on 2025-01-15).}


%% ═══════════════════════════════════════════════════════════════
\section{D2: Compositional Extraction}
\label{sec:d2}
%% ═══════════════════════════════════════════════════════════════

\begin{definition}[Desideratum D2 --- Compositional Extraction]
\label{def:d2}
In the primary regime, querying a revised state should decompose into
queries of the components:
\[
  \mathrm{extract}(\mathrm{revise}(W_1, W_2),\, q)
  = \mathrm{combine}(\mathrm{extract}(W_1, q),\;
                       \mathrm{extract}(W_2, q)).
\]
That is, extraction should be a homomorphism from the state monoid to
the value monoid.
\end{definition}

\noindent
This is the algebraic heart of the additive regime.  It says that
evidence from independent sources can be combined without interference:
the result of revising and then querying equals the result of querying
each source separately and then combining the answers.

\paragraph{Why this matters.}
Compositional extraction is what makes the profile-hom equivalence
(Theorem~\ref{thm:universal-property}) possible: an additive world
model \emph{is} a monoid homomorphism from states to answer profiles.
Without compositionality, there is no universal property, no clean
categorical picture, and no guarantee that compiled inference shortcuts
agree with full extraction.

\paragraph{The foundational justification.}
D2 is not an arbitrary algebraic choice.  De~Finetti's representation
theorem for exchangeable sequences shows that, under exchangeability,
sufficient statistics add: $(n_1^+, n_1^-) + (n_2^+, n_2^-) =
(n_1^+ + n_2^+, n_1^- + n_2^-)$.  The Solomonoff--de~Finetti
foundation (Chapter~\ref{ch:world-models}, \S``The Solomonoff--De
Finetti Foundation'') establishes the categorical necessity chain:
exchangeability $\to$ de~Finetti mixture $\to$ Kyburg flattening $\to$
Giry bind $\to$ evidence sufficiency.  D2 is the algebraic shadow of
this necessity.

\paragraph{What D2 rules out.}
Any system whose query answers depend on the \emph{order} of evidence
arrival (in the independent-sources case) violates D2.  Classical
Bayesian updating satisfies D2 at the level of sufficient statistics
but violates it at the level of posterior probabilities (normalization
destroys additivity).  This is exactly why the calculus operates on
carriers (sufficient statistics), not on probabilities.

\paragraph{What D2 does not require.}
D2 is stated for the ``primary regime'' (the additive regime).  It is
\emph{not} required for all revision modes.  The overlap-aware,
trust-gated, and coalition regimes of Chapter~\ref{ch:world-models}
explicitly violate D2, and this is expected: when sources share
information, or when trust gates revision, naive composition is wrong.
D2 identifies the regime where composition is exact; the other regimes
are frontier directions where D2's algebraic convenience is traded for
richer structure.

\paragraph{Discharge.}
D2 is discharged by the \code{AdditiveWorldModel} axiom
\code{extract\_add} (Chapter~\ref{ch:world-models},
\S\ref{sec:additive-regime}) and the profile-hom equivalence
(Theorem~\ref{thm:universal-property}).

\paragraph{Independence witness.}
A system that carries states and supports forgetting (D1, D6) but
uses a nonlinear extraction---e.g.,
$\mathrm{extract}(W_1 + W_2, q) = \max(\mathrm{extract}(W_1, q),
\mathrm{extract}(W_2, q))$---satisfies D1, D3--D8 but violates D2.

\maplecourt{%
Apartment~3B receives two independent sensor reports about the pipe:
$(3, 1)$ from the north sensor and $(2, 0)$ from the south sensor.
D2 says the combined evidence is $(5, 1)$, the same whether we
combine first and then query, or query each and combine the answers.
The order of arrival does not matter; the combination is commutative
and associative.  This is why the apartment can buffer overnight
readings and batch-update at dawn (the sleep-consolidation theorem,
Theorem~\ref{thm:sleep-consolidation}).}


%% ═══════════════════════════════════════════════════════════════
\section{D3: Generic Evidence Carriers}
\label{sec:d3}
%% ═══════════════════════════════════════════════════════════════

\begin{definition}[Desideratum D3 --- Generic Carriers]
\label{def:d3}
The value type $V$ in which extracted answers live should be a
\emph{parameter}, not a fixed datatype.  Different domains (binary
yes/no, multi-valued categories, continuous measurements) should
instantiate the same algebraic discipline with different carriers,
and the calculus should not change across them.
\end{definition}

\noindent
The motivation is domain coverage.  A framework that handles only
binary queries is too narrow for practical applications.  Elevator
health is three-valued (normal, slow, faulty).  Hallway humidity is
continuous.  Gene expression levels are real-valued vectors.  Each
domain has its own natural sufficient statistics, and a good calculus
should accommodate all of them under one roof.

\paragraph{The algebraic discipline.}
D3 does not ask for $V$ to be completely arbitrary.  It asks for $V$ to
be an \code{AddCommMonoid}: a commutative monoid under addition with a
zero element.  This is the minimal structure needed for D2
(compositional extraction) to make sense.  Each carrier instantiates
this structure with its own sufficient statistics:

\begin{center}
\begin{tabular}{lll}
\toprule
\textbf{Carrier} & \textbf{$V$} & \textbf{AddCommMonoid operation} \\
\midrule
Binary & $(n^+, n^-) \in \mathbb{R}_{\ge 0}^2$
  & componentwise addition \\
Dirichlet & $(\alpha_1, \ldots, \alpha_k) \in \mathbb{R}_{\ge 0}^k$
  & componentwise addition \\
Normal-Gamma & $(n, \sum x_i, \sum x_i^2) \in \mathbb{R}_{\ge 0}^3$
  & componentwise addition \\
Amplitude & $z \in \mathbb{C}$
  & complex addition \\
\bottomrule
\end{tabular}
\end{center}

\noindent
The same five algebraic rules (Chapter~\ref{ch:world-models},
\S\ref{sec:algebraic-core}) hold for every row of this table.  The
calculus changes only in the choice of carrier; the algebra stays
the same.

\paragraph{What D3 rules out.}
Any framework that hard-codes a single evidence type (e.g., binary
strength-confidence pairs) violates D3.  PLN v0.9 uses binary evidence
exclusively; the world-model calculus treats binary evidence as one
instance among several (Chapter~\ref{ch:evidence}).

\paragraph{Discharge.}
D3 is discharged by the parametric definition
$\code{WorldModel (State Query V : Type*)}$ with $V$ constrained
only by \code{[AddCommMonoid V]} in the additive regime.  The three
formalized carriers (binary, Dirichlet, Normal-Gamma) are concrete
instances (Chapter~\ref{ch:evidence}).

\paragraph{Independence witness.}
PLN v0.9 (with binary evidence only) satisfies D1, D2, D4--D8 but
violates D3: the carrier is not generic.

\maplecourt{%
Maple Court uses all three formalized carriers simultaneously:
binary $(n^+, n^-)$ for leak alarms and occupancy; Dirichlet
$(\alpha_1, \alpha_2, \alpha_3)$ for the three-valued elevator
health and laundry status; Normal-Gamma $(n, \sum x_i, \sum x_i^2)$
for continuous humidity readings.  D3 ensures that the same revision
law, the same extraction interface, and the same forgetting mechanism
work across all three---only the carrier changes.}


%% ═══════════════════════════════════════════════════════════════
\section{D4: Evidence over Views}
\label{sec:d4}
%% ═══════════════════════════════════════════════════════════════

\begin{definition}[Desideratum D4 --- Evidence over Views]
\label{def:d4}
Operational readouts---strength, confidence, interval bounds, posterior
means---are \emph{derived projections} from richer evidence carriers.
They must never be treated as the primary semantic objects.  Evidence
is what you carry; views are what you display.
\end{definition}

\noindent
This is the desideratum that most sharply distinguishes the world-model
calculus from its predecessors.  Classical probability theory operates
on probabilities.  PLN v0.9 operates on strength-confidence pairs.
NARS operates on frequency-confidence pairs.  All of these treat the
projected value as the thing being stored and manipulated.  D4 says:
store the evidence, not the projection.

\paragraph{Why views are not proof objects.}
A strength value $s = 0.83$ tells you nothing about how much evidence
produced it.  The evidence pairs $(5, 1)$ and $(500, 100)$ both have
strength $\approx 0.83$, but the second is far more informative.
More importantly, views are not additive
(Proposition~\ref{prop:overview-views-not-semantic-core}): the
strength of combined evidence is not the sum of individual strengths.
Any system that propagates views through inference chains accumulates
artificial precision---the five chaining-limitation theorems of
Chapter~\ref{ch:error-transport} quantify this loss precisely.

\paragraph{The categorical reading.}
In categorical terms, evidence carriers are objects in the category
of additive commutative monoids; views are morphisms from carriers to
simpler monoids (or non-monoids).  The key fact is that most views
are \emph{not} monoid homomorphisms: they do not preserve the
additive structure.  This is why views cannot serve as the semantic
layer---they break the algebraic structure that D2 provides.

\paragraph{What D4 rules out.}
Any system that stores only $s$ and $c$ (strength and confidence),
discarding the underlying counts, violates D4.  Such a system cannot
support forgetting (you cannot retract part of a scalar), cannot
distinguish high-confidence low-evidence from high-confidence
high-evidence, and cannot propagate full posteriors through inference
chains.

\paragraph{Discharge.}
D4 is discharged by the architecture of the calculus: extraction
returns values in the carrier $V$ (the evidence layer), and views
are defined as separate functions on $V$ (Chapter~\ref{ch:views}).
The strength non-monotonicity theorem
(Theorem~\ref{thm:strength-not-monotone}) proves that views are not
order-preserving, making the separation essential rather than
conventional.

\paragraph{Independence witness.}
A system that stores full evidence carriers but has no extraction
homomorphism (violating D2) still satisfies D4.  Conversely, a system
with compositional extraction into a carrier isomorphic to its views
(e.g., $V = \mathbb{R}$ with extraction returning strength directly)
satisfies D2 but violates D4.


%% ═══════════════════════════════════════════════════════════════
\section{D5: Explicit Side Conditions}
\label{sec:d5}
%% ═══════════════════════════════════════════════════════════════

\begin{definition}[Desideratum D5 --- Explicit Side Conditions]
\label{def:d5}
Every compiled inference shortcut must carry its structural
assumptions as a typed, machine-checkable predicate $\Sigma$.  When
$\Sigma$ holds, the shortcut is \emph{exact}: it agrees with full
extraction.  When $\Sigma$ fails, the shortcut is not licensed.
There are no hidden assumptions.
\end{definition}

\noindent
This desideratum addresses a deep problem in practical uncertain
reasoning: rules that work under unstated assumptions.  The PLN
deduction formula is correct---but only under screening-off and
positivity ($0 < s_B < 1$).  Naive Bayes classification is
exact---but only under the conditional independence assumption encoded
in a fork BN.  D5 says: state the assumption.  Prove the rule.  Check
the assumption at runtime.

\paragraph{The $\Sigma$-discipline.}
In the calculus, every compiled rewrite has the form:
\[
  \frac{\Sigma \vdash \mathrm{sideCondition} \quad
        \Sigma; \Gamma \vdash \mathrm{premises}}
       {\Sigma; \Gamma \vdash \mathrm{conclusion}}
\]
The side condition $\Sigma$ appears above the line.  It is not
optional.  It is a typed proposition that can be proved, checked by a
decision procedure, or required as a hypothesis.

\paragraph{What D5 rules out.}
Any system that applies inference rules unconditionally---without
checking or even stating their side conditions---violates D5.  PLN
v0.9 states rules as universal formulas (``deduction always gives
$s_{AC} = \ldots$''); the world-model calculus states them as
conditional theorems (``if screening-off and positivity, then
$s_{AC} = \ldots$'').  The no-go theorem
(Chapter~\ref{ch:no-go}) shows that unconditional application is
unsound in general: two world-model states can agree on all local
per-link summaries while differing on the target query.

\paragraph{Discharge.}
D5 is discharged by the $\Sigma$-guarded rewrite architecture of
Chapter~\ref{ch:sigma-rewrites} and the BN exactness theorems of
Chapter~\ref{ch:bn-models}.  Each compiled rule in the catalog
(Chapter~\ref{ch:sigma-rewrites}, \S``The Compiled Rule Catalog'')
has an explicit $\Sigma$ column.

\paragraph{Independence witness.}
A purely extensional world model that stores full joint distributions
and answers queries by marginalization satisfies D1--D4, D6--D8 but
has no compiled rules and therefore trivially satisfies (or vacuously
satisfies) D5.  The interesting violation is a system \emph{with}
compiled rules but \emph{without} explicit side conditions---this is
essentially PLN v0.9, which violates D5 while approximately
satisfying the others.

\maplecourt{%
The compiled chain rule for
$\mathrm{PipeLeak} \to \mathrm{WallHumidity} \to \mathrm{MoldRisk}$
carries the side condition
$\Sigma = \{\mathrm{ScreenOff}(\mathrm{PipeLeak},
\mathrm{WallHumidity}, \mathrm{MoldRisk}),\;
0 < s_{\mathrm{Humidity}} < 1\}$.
When this $\Sigma$ holds---which it does for Maple Court's BN
structure, where the DAG guarantees d-separation---the compiled
formula gives the exact posterior.  When an additional direct
leak-to-mold pathway is added (breaking screening-off), $\Sigma$
fails and the rule is not licensed.  D5 ensures that this failure is
detected, not silently ignored.}


%% ═══════════════════════════════════════════════════════════════
\section{D6: Retractability}
\label{sec:d6}
%% ═══════════════════════════════════════════════════════════════

\begin{definition}[Desideratum D6 --- Retractability]
\label{def:d6}
It must be possible to \emph{retract} evidence within a declared scope
while provably preserving all answers outside that scope.
Formally: for every scope~$S$ and state~$W$, there exists a
forgetting operation $\mathrm{forget}(S, W)$ such that
\[
  \forall q \notin S,\quad
  \mathrm{extract}(\mathrm{forget}(S, W), q) = \mathrm{extract}(W, q).
\]
\end{definition}

\noindent
Retractability is what separates stateful reasoning from stateless
truth-value manipulation.  If a sensor is found to be miscalibrated,
its readings must be retractable.  If a privacy constraint requires
apartment-level data to expire after 24~hours, the system must forget
that data without corrupting building-level answers.  If a correction
arrives (``observation~3 was a false alarm''), the system must remove
that observation's contribution while preserving everything else.

\paragraph{Why retractability is not deletion.}
Deleting evidence is easy: set the state to zero.  Retractability is
hard because it requires \emph{surgical} removal: forget the evidence
within scope~$S$ while proving that answers outside~$S$ are unchanged.
This requires the state to carry enough structure to distinguish what
came from where---which is precisely what provenance provides
(Chapter~\ref{ch:state-ops}, \S\ref{sec:provenance}).

\paragraph{What D6 rules out.}
Any system that cannot undo a revision---either because it stores only
aggregate summaries (violating D1) or because it lacks provenance
tracking---violates D6.  Classical Bayesian updating can, in principle,
undo a conjugate update by subtracting the likelihood contribution,
but this requires knowing which data point to subtract.  Without
provenance, the data point is lost in the aggregate.

\paragraph{Discharge.}
D6 is discharged by the \code{ForgettingLayer} structure and the
outside-invariant law of Chapter~\ref{ch:state-ops},
\S\ref{sec:forgetting}.

\paragraph{Independence witness.}
An append-only evidence store (states grow monotonically, no
retraction) satisfies D1--D5, D7--D8 but violates D6.

\maplecourt{%
Three hallway humidity alerts arrive, then sensor~3B-north is found
to be miscalibrated.  If the system stores only the aggregate
$(5, 1)$, it cannot retract sensor~3B-north's contribution without
knowing how much of the $(5, 1)$ came from that sensor.  With
provenance, the state records
$\{(\text{3B-north}, (3, 1)), (\text{3B-south}, (2, 0))\}$, and
forgetting scope $S = \text{3B-north}$ yields $(2, 0)$.  The
building-level humidity answer (which does not depend on 3B's pipe
sensor) is unchanged.}


%% ═══════════════════════════════════════════════════════════════
\section{D7: Evidence Conservation}
\label{sec:d7}
%% ═══════════════════════════════════════════════════════════════

\begin{definition}[Desideratum D7 --- Evidence Conservation]
\label{def:d7}
Evidence can be neither \emph{fabricated} (inference cannot create
support that was not present in the premises) nor \emph{overcounted}
(the same observation cannot contribute twice to the same query).
Formally, three conservation laws hold under exact-inverse forgetting:
\begin{enumerate}
\item \textbf{Anti-hallucination}: a revision $\Delta$ within scope
  $S$ contributes zero evidence to queries outside~$S$.
\item \textbf{Noether conservation}: the revised state agrees with the
  base state on queries outside the revision's scope.
\item \textbf{Zero leakage}: the observation count of $\Delta$ at
  queries outside its scope is zero.
\end{enumerate}
\end{definition}

\noindent
D7 is the evidence-theoretic analogue of conservation laws in physics.
Just as energy cannot be created or destroyed, evidence cannot be
fabricated or lost.  The ``anti-hallucination'' property is
particularly important in the context of large language models and
generative AI: a system satisfying D7 cannot output a confident
assertion unless the evidence to support it actually exists in the
state.

\paragraph{Two faces of conservation.}
The calculus establishes conservation from two perspectives:
\begin{itemize}
\item \textbf{State-level} (Chapter~\ref{ch:state-ops}): the
  conservation theorems (Theorem~\ref{thm:conservation}) show that
  scoped forgetting preserves evidence outside the scope.
\item \textbf{Inference-path level}
  (Goertzel~\cite{GoertzelEvidenceConservation2026}): a discrete
  quantale Noether theorem showing that reinforcement is constant
  along geodesic inference paths, a hallucination bound on conclusion
  strength, and an evidence monotonicity theorem (a quantale
  data-processing inequality).
\end{itemize}
Together, the two perspectives show that evidence integrity holds
whether you look at state retraction or at inference chaining.

\paragraph{What D7 rules out.}
Any system where chaining inference rules can amplify evidence beyond
what the premises contain violates D7.  This includes systems where
deduction applied twice produces stronger conclusions than the
original evidence warrants.  The stamp-disjointness condition in the
revision rule (Chapter~\ref{ch:sigma-rewrites}) is one mechanism that
enforces D7: stamps prevent the same observation from being counted
twice.

\paragraph{Discharge.}
D7 is discharged by the evidence conservation pack
(\code{EvidenceConservationPack} in
Chapter~\ref{ch:state-ops}, \S``Conservation Theorems'') and the
exact-inverse no-go theorem (if evidence leaked, forgetting cannot
exactly undo it).

\paragraph{Independence witness.}
A system with scoped forgetting (D6) but no conservation guarantee
(e.g., a sloppy retraction that removes too much or too little)
satisfies D6 but violates D7.  Conversely, a system that never forgets
(violating D6) trivially satisfies D7 by never modifying the state.

\maplecourt{%
Suppose the building model infers mold risk from pipe-leak evidence
via a chain of compiled rules.  D7 guarantees that the mold-risk
evidence in the conclusion cannot exceed the pipe-leak evidence in
the premises.  If only $(5, 1)$ of leak evidence exists, the
mold-risk conclusion has at most $(5, 1)$ worth of evidential
support (modulated by the chain's conditional probabilities).  The
anti-hallucination property says: if the pipe-leak evidence is
entirely within scope ``apartment 3B,'' then the mold-risk answer
for apartment 5C is unaffected.}


%% ═══════════════════════════════════════════════════════════════
\section{D8: Honest Boundaries}
\label{sec:d8}
%% ═══════════════════════════════════════════════════════════════

\begin{definition}[Desideratum D8 --- Honest Boundaries]
\label{def:d8}
The calculus must state what it \emph{cannot} do as precisely as what
it can do.  Impossibility results---no-go theorems, coverage collapse
bounds, non-uniqueness demonstrations---are first-class contributions,
not embarrassments.  Every boundary theorem has the same formal status
as every positive theorem.
\end{definition}

\noindent
D8 is a meta-desideratum: it constrains how the calculus presents
itself.  A framework that claims universality without stating its
limits is either trivial or dishonest.  The world-model calculus earns
its credibility by proving both what works and what does not.

\paragraph{The boundary catalog.}
The book contains a substantial family of boundary results:
\begin{itemize}
\item \textbf{No-go for local link completeness}
  (Chapter~\ref{ch:no-go}): per-link summaries cannot reconstruct
  the joint.
\item \textbf{Five chaining limitations}
  (Chapter~\ref{ch:error-transport}): variance accumulation,
  sensitivity amplification, coverage collapse, topology-CPT gap,
  and no universal error bound.
\item \textbf{Inclusion-exclusion identifiability no-go}
  (Chapter~\ref{ch:scale-control}): the first two
  inclusion-exclusion terms do not determine the union cardinality.
\item \textbf{Trail-free non-stabilization}
  (Chapter~\ref{ch:scale-control}): re-deriving conclusions without
  trails can oscillate.
\item \textbf{Pooling non-uniqueness}
  (Chapter~\ref{ch:scale-control}): the four pooling axioms do not
  pin down additive fusion.
\item \textbf{No binary mixed policy}
  (Chapter~\ref{ch:intensional}): the classical two-channel
  intensional mixture is refuted.
\item \textbf{Exact-inverse no-go}
  (Chapter~\ref{ch:state-ops}): if evidence leaked, forgetting
  cannot exactly undo it.
\item \textbf{Strength non-monotonicity}
  (Chapter~\ref{ch:views}): more evidence can yield lower strength.
\end{itemize}

\noindent
Each boundary theorem carries a \claimBoundary{} label, making it
visually and taxonomically distinct from \claimCore{} and
\claimConditional{} results.

\paragraph{What D8 rules out.}
Any framework that presents only positive results---``our system can
do X, Y, Z''---without stating where it fails violates D8.  This
includes marketing-style presentations of AI systems (``universal
reasoning engine'') that omit computational complexity, soundness
gaps, or regime limitations.

\paragraph{Discharge.}
D8 is discharged by the boundary catalog above and by the claim
taxonomy that labels every theorem as \claimCore{}, \claimConditional{},
or \claimBoundary{}.

\paragraph{Independence witness.}
A framework with correct positive results but no proved impossibility
theorems violates D8 while potentially satisfying D1--D7.


%% ═══════════════════════════════════════════════════════════════
\section{The Desiderata as a Whole}
\label{sec:desiderata-summary}
%% ═══════════════════════════════════════════════════════════════

The eight desiderata form a layered specification:

\begin{center}
\begin{tabular}{clll}
\toprule
\textbf{ID} & \textbf{Name} & \textbf{Layer}
  & \textbf{Primary discharge} \\
\midrule
D1 & State Primacy & Semantic kernel
  & Ch~\ref{ch:world-models}: \code{WorldModel} \\
D2 & Compositional Extraction & Semantic kernel
  & Ch~\ref{ch:world-models}: profile-hom \\
D3 & Generic Carriers & Evidence architecture
  & Ch~\ref{ch:evidence}: three carriers \\
D4 & Evidence over Views & Evidence architecture
  & Ch~\ref{ch:views}: lossy projections \\
D5 & Explicit Side Conditions & Inference discipline
  & Ch~\ref{ch:sigma-rewrites}: $\Sigma$-guards \\
D6 & Retractability & State management
  & Ch~\ref{ch:state-ops}: forgetting layer \\
D7 & Evidence Conservation & State management
  & Ch~\ref{ch:state-ops}: conservation pack \\
D8 & Honest Boundaries & Self-knowledge
  & Ch~\ref{ch:no-go}: no-go theorems \\
\bottomrule
\end{tabular}
\end{center}

\paragraph{Necessity.}
Are all eight necessary?  The independence witnesses show that dropping
any one produces a system that, while coherent, fails in a specific and
important way: without D1, you cannot retract; without D2, you cannot
compose; without D3, you are limited to one domain; without D4, you
confuse projections with semantics; without D5, your shortcuts are
unsound; without D6, your state is append-only; without D7, your
inference can hallucinate; without D8, you do not know your own limits.

\paragraph{Sufficiency.}
Are the eight jointly sufficient for ``responsible uncertain
reasoning''?  Not in full generality.  At least three additional
properties are needed for a complete system but are not yet fully
developed in this book:
\begin{itemize}
\item \textbf{Computational tractability}: the calculus says
  \emph{what} the correct answer is but does not bound the
  \emph{cost} of computing it.
\item \textbf{Non-additive regimes}: D2 holds only in the additive
  regime; overlap-aware and trust-gated revision need their own
  algebraic discipline.
\item \textbf{Measure-theoretic grounding}: the current development
  works with finitary carriers; the terminal measure-valued world
  model is an open program.
\end{itemize}
These are recorded as frontier directions (Chapter~\ref{ch:future}),
not as current desiderata.  The eight listed here are the contract
that the present book fulfills.

\paragraph{The PLN heritage.}
PLN v0.9 satisfies D1 and D2 in spirit, partially satisfies D7 (via
evidence stamps), but does not satisfy D3 (single carrier), D4 (scalar
propagation), D5 (no explicit $\Sigma$), D6 (no formal forgetting),
or D8 (no proved no-go theorems).  The world-model calculus inherits
PLN's conceptual vocabulary and evidence-based philosophy while adding
the formal machinery that PLN lacked.  The desiderata are the precise
statement of what that machinery is for.

\begin{figure}[htbp]
\centering
\begin{tikzpicture}[
  >=Latex,
  every node/.style={font=\footnotesize},
  dbox/.style={draw, rounded corners=4pt, minimum width=32mm,
              minimum height=10mm, align=center},
]
  % Kernel layer
  \node[dbox, fill=blue!8] (d1) at (0, 4) {D1: State Primacy};
  \node[dbox, fill=blue!8] (d2) at (5, 4) {D2: Composition};

  % Evidence layer
  \node[dbox, fill=green!8] (d3) at (0, 2.5) {D3: Generic Carriers};
  \node[dbox, fill=green!8] (d4) at (5, 2.5) {D4: Evidence $>$ Views};

  % Inference layer
  \node[dbox, fill=yellow!12] (d5) at (0, 1) {D5: Explicit $\Sigma$};
  \node[dbox, fill=yellow!12] (d6) at (5, 1) {D6: Retractability};

  % Meta layer
  \node[dbox, fill=orange!10] (d7) at (0, -0.5) {D7: Conservation};
  \node[dbox, fill=orange!10] (d8) at (5, -0.5) {D8: Honest Boundaries};

  % Layer labels
  \node[font=\scriptsize, text=blue!60!black, rotate=90] at (-2.8, 3.25) {kernel};
  \node[font=\scriptsize, text=green!50!black, rotate=90] at (-2.8, 2.5) {evidence};
  \node[font=\scriptsize, text=yellow!60!black, rotate=90] at (-2.8, 1) {inference};
  \node[font=\scriptsize, text=orange!60!black, rotate=90] at (-2.8, -0.5) {integrity};

  % Dependency arrows
  \draw[->, thick, gray!50] (d1) -- (d6);
  \draw[->, thick, gray!50] (d2) -- (d5);
  \draw[->, thick, gray!50] (d3) -- (d4);
  \draw[->, thick, gray!50] (d6) -- (d7);
  \draw[->, thick, gray!50] (d5) -- (d8);

  % Annotation
  \node[font=\scriptsize, text=gray, text width=40mm, align=center]
    at (2.5, -1.8) {Arrows show enablement:\\
    D1 enables D6 (states enable retraction);\\
    D2 enables D5 (composition enables shortcuts);\\
    D6 enables D7 (retraction enables conservation).};
\end{tikzpicture}
\caption{The eight desiderata organized by layer.  Arrows show
enablement relationships: D1 (state primacy) enables D6
(retractability); D2 (compositional extraction) enables D5 (explicit
side conditions for compiled shortcuts); D6 (retractability) enables
D7 (conservation).}
\label{fig:desiderata-layers}
\end{figure}

\maplecourt{%
Maple Court's building model satisfies all eight desiderata.
It carries states (D1), combines sensor reports additively (D2),
handles binary, Dirichlet, and Normal-Gamma evidence (D3), derives
strength and confidence as views (D4), gates compiled chain rules on
BN screening-off (D5), retracts miscalibrated sensor data (D6),
conserves evidence under retraction (D7), and reports where local
chaining fails (D8).  The desiderata are not abstract stipulations
imposed from outside---they are the properties that make a building
automation system trustworthy.}
